\section{Discussion}
The fallback recovered cases for which the process claim was missing but the original endpoints, authentication claim, and network claim agreed. It added 18 qualifying groups anywhere in the queue, but only ten additional top-20 recalls. The eight-incident gap is a reminder to look at ranking after reconstruction. Perfect observed grouped-pair precision on this sample also says nothing about whether real upstream sources would be independent.

Two details drove the evidence rule: a copied record retains its upstream root, and a claim is usable only after it arrives, regardless of its event timestamp. Separate development stress tests lost recovery when evidence was delayed, absent, corrupted, or contradictory. A fourth independent source was unavailable, so we did not invent one to fill the gap. Ranking-only trials were also inconclusive. For example, a pairwise-trained score moved development recall from 263/300 to 262/300 with identical groups; its training objective did not match the evaluator's top-20 target.

There is a cost to the strict rule. It abstained on 35 genuine proposed relationships in the holdout even though no false grouped pairs were observed. Whether an analyst would prefer this tradeoff depends on missed-incident costs, telemetry quality, and the review process. We did not measure any of those operational quantities.
